\section{Introduction}
\label{sec:intro}

Source code is now the primary attack surface of most organisations. A
single hard-coded credential, an unsanitised path into a query builder,
or an unsafe deserialisation call can compromise an entire deployment,
and the categories of weakness responsible have remained remarkably
stable for two decades~\cite{owasp2021,mitrecwe}. Two families of tools
have emerged to help developers find such defects before release. Static
application security testing (SAST) analyses code without executing it,
using pattern matching, data-flow, and taint tracking to flag suspicious
constructs~\cite{chess2004,ayewah2008,yamaguchi2014}. More recently,
large language models (LLMs) trained on code have been applied to review,
vulnerability detection, and repair, offering a form of semantic
reasoning that rule-based analysers lack~\cite{chen2021codex,pearce2023,
khoury2023}.

Despite their complementary strengths, both families share a deployment
assumption that is increasingly at odds with how sensitive code is
governed: they run in the cloud. Commercial SAST platforms upload the
repository to a hosted analysis backend, and virtually every capable
code-reasoning LLM is accessed through a remote API. For teams working on
proprietary algorithms, unreleased products, defence or medical software,
or anything subject to data-residency rules, transmitting the source
tree off the premises is either prohibited outright or requires a
contractual and architectural effort that discourages routine use. The
same barrier blocks adoption in genuinely air-gapped settings. The
practical consequence is that the developers who arguably need rigorous
auditing most are the least able to use the current generation of
assistive tools.

A second, well-documented problem is that developers frequently abandon
static analysers because the tools do not fit their workflow: results
arrive out of context, the signal-to-noise ratio is poor, and the
interface makes triage tedious~\cite{johnson2013,nachtigall2022}. A tool
that solves the privacy problem but reproduces the usability problem will
not be adopted either.

This paper describes \cs, a desktop application that addresses both
concerns. \cs performs a complete security and structural audit of a
repository \emph{entirely on the developer's own machine}. It combines
three analysis modalities---pattern-based static detection, on-device
LLM semantic review, and containerised dynamic execution---behind a
single graphical workflow, and it aggregates their output, together with
classical software metrics, into one interpretable score. The LLM
component runs a quantised \mbox{Llama~3.2} model inside a local
container via Ollama~\cite{ollama,touvron2023llama2}; the dynamic
component builds and runs the project inside a resource-limited Docker
sandbox~\cite{merkel2014docker}. No stage of the pipeline opens an
outbound network connection, and the application exposes an explicit
offline indicator so that this property is visible to the user.

\smallskip
\noindent\textbf{Contributions.} This paper makes the following
contributions:
\begin{itemize}
  \item A \emph{local-first architecture} that co-locates pattern-based
        SAST, an on-device code-reasoning LLM, a cyclomatic-complexity
        engine, and a container sandbox in a single desktop process, with
        a hard no-egress boundary (\cref{sec:method}).
  \item A \emph{finding-fusion and composite-scoring} method that
        reconciles rule-based and model-based findings and folds them,
        with structural metrics, into a single \phindex\ and a set of
        interpretable risk bands (\cref{sec:fusion,sec:risk}).
  \item A \emph{task-oriented interface} of twelve coordinated screens
        that let a developer configure an audit, watch it run, and triage
        results in the context of the offending source
        (\cref{sec:interface}).
  \item A \emph{worked case study} on a representative
        retrieval-augmented-generation (RAG) service that exercises the
        weakness classes now specific to LLM-integrated
        applications~\cite{owaspllm2025,greshake2023}, and a discussion of
        where local inference helps and where it costs recall
        (\cref{sec:results,sec:discussion}).
\end{itemize}

\smallskip
\noindent The remainder of the paper is organised as follows.
\Cref{sec:related} surveys static and dynamic analysis, software metrics,
and the use of language models for code and security.
\Cref{sec:method} presents the architecture, the analysis pipeline, and
the scoring model. \Cref{sec:interface} walks through the product
interface. \Cref{sec:results} reports the case study. \Cref{sec:discussion}
discusses trade-offs and threats to validity, and \Cref{sec:conclusion}
concludes.
